% =================================================
% CHAPITRE 4 : SYSTEMES LINEAIRES
% =================================================

\chapter{Systèmes linéaires}
\label{chap-systemes}

\section{Définition}

\begin{definition}[Système linéaire]
Un \emph{système linéaire} de $n$ équations à $n$
inconnues $x_1, \dots, x_n$ s'écrit
\[
    \left\{
    \begin{array}{ccccccc}
        a_{11} x_1 & + & \cdots & + & a_{1n} x_n & = & b_1 \\
        \vdots     &   &        &   & \vdots     &   & \vdots \\
        a_{n1} x_1 & + & \cdots & + & a_{nn} x_n & = & b_n
    \end{array}
    \right.
\]
Avec le produit d'une matrice par un vecteur
(\cref{def-matrice-vecteur}), il s'écrit simplement
\[
    A\vect{x} = \vect{b},
    \qquad
    A \in \R^{n \times n},\
    \vect{x} \in \R^n,\
    \vect{b} \in \R^n.
\]
$A$ est la \emph{matrice du système}, $\vect{b}$ le
\emph{second membre} et $\vect{x}$ le vecteur des
inconnues.
\end{definition}

\begin{definition}[Matrice augmentée]
Pour résoudre le système à la main, on écrit $A$ et
$\vect{b}$ côte à côte dans la \emph{matrice augmentée}
$[A \mid \vect{b}] \in \R^{n \times (n+1)}$ :
\[
    \left(
    \begin{array}{ccc|c}
        a_{11} & \cdots & a_{1n} & b_1 \\
        \vdots &        & \vdots & \vdots \\
        a_{n1} & \cdots & a_{nn} & b_n
    \end{array}
    \right).
\]
\end{definition}

\begin{definition}[Matrice inversible]
Une matrice carrée $A \in \R^{n \times n}$ est
\emph{inversible} s'il existe une matrice
$A^{-1} \in \R^{n \times n}$ telle que
\[
    A A^{-1} = A^{-1} A = I_n.
\]
Dans ce cas, le système $A\vect{x} = \vect{b}$ a une
solution unique pour tout $\vect{b}$ :
$\vect{x} = A^{-1}\vect{b}$. Une matrice qui n'est
pas inversible est dite \emph{singulière}.
\end{definition}

\begin{importantbox}
On ne calcule presque jamais $A^{-1}$ pour résoudre
$A\vect{x} = \vect{b}$ : la méthode de Gauss donne
$\vect{x}$ directement, plus vite et avec moins
d'erreurs d'arrondi.
\end{importantbox}

\section{Le pivot de Gauss}

\subsection{Opérations élémentaires}

\begin{proposition}
Les trois opérations suivantes sur les lignes de la
matrice augmentée ne changent pas les solutions du
système :
\begin{enumerate}
    \item échanger deux lignes :
          $L_i \leftrightarrow L_j$ ;
    \item multiplier une ligne par un réel non nul :
          $L_i \leftarrow \lambda L_i$, avec
          $\lambda \neq 0$ ;
    \item ajouter à une ligne un multiple d'une autre
          ligne : $L_i \leftarrow L_i - \ell\, L_j$,
          avec $i \neq j$.
\end{enumerate}
\end{proposition}

\subsection{Élimination}
\label{sec-elimination}

\begin{definition}[Matrice triangulaire]
Une matrice carrée $U$ est \emph{triangulaire
supérieure} si tous ses coefficients sous la
diagonale sont nuls : $u_{ij} = 0$ pour $i > j$. Elle
est \emph{triangulaire inférieure} si tous ses
coefficients au-dessus de la diagonale sont nuls.
\end{definition}

L'idée de Gauss : utiliser l'opération 3 pour
transformer $A\vect{x} = \vect{b}$ en un système
\emph{triangulaire supérieur} $U\vect{x} = \vect{c}$,
qui a les mêmes solutions et qui se résout facilement.

\begin{methodbox}
Pour chaque colonne $k = 1, \dots, n - 1$ :
\begin{enumerate}
    \item le \emph{pivot} est le coefficient
          diagonal $a_{kk}$ ;
    \item pour chaque ligne $i > k$ sous le pivot,
          calculer le \emph{multiplicateur}
          $\ell_{ik} = \dfrac{a_{ik}}{a_{kk}}$ ;
    \item faire $L_i \leftarrow L_i - \ell_{ik}\, L_k$,
          sur toute la ligne de la matrice augmentée
          (second membre compris). Le coefficient
          $a_{ik}$ devient $0$.
\end{enumerate}
À la fin, tous les coefficients sous la diagonale
sont nuls.
\end{methodbox}

\subsection{Remontée}

\begin{methodbox}
Le système $U\vect{x} = \vect{c}$ se résout de
\textbf{bas en haut}. La dernière ligne ne contient
qu'une inconnue, puis chaque ligne n'en ajoute
qu'une :
\[
    x_n = \frac{c_n}{u_{nn}},
    \qquad
    x_i = \frac{1}{u_{ii}}
    \left( c_i - \sum_{j = i+1}^{n} u_{ij}\, x_j \right)
    \quad \text{pour } i = n - 1, \dots, 1.
\]
\end{methodbox}

\begin{importantbox}
En C++, la boucle de la remontée descend de $n - 1$
à $0$ inclus. Avec un indice de type
\cpp|std::size_t| (non signé), la condition
\cpp|i >= 0| est toujours vraie : la boucle ne
s'arrête jamais. Utiliser un \cpp|int|.
\end{importantbox}

\subsection{Le pivot partiel}

La méthode de la \cref{sec-elimination} divise par
le pivot $a_{kk}$.
Elle échoue si ce pivot est nul, alors même que le
système a une solution unique.

\begin{definition}[Pivot partiel]
Avant de traiter la colonne $k$, on cherche, parmi les
lignes $k$ à $n$, celle dont le coefficient de la
colonne $k$ a la plus grande \textbf{valeur absolue} :
\[
    p = \operatorname*{arg\,max}_{k \le i \le n} \abs{a_{ik}}.
\]
On échange les lignes $k$ et $p$ (second membre
compris), puis on élimine comme avant.
\end{definition}

\begin{remark}
Le pivot partiel ne sert pas seulement à éviter la
division par $0$. Un pivot \emph{très petit} donne
des multiplicateurs $\ell_{ik}$ très grands, qui
amplifient les erreurs d'arrondi. Choisir le plus
grand pivot garantit $\abs{\ell_{ik}} \le 1$.
\end{remark}

\begin{definition}[Pivot numériquement nul]
En \cpp|double|, un calcul qui devrait donner $0$
donne souvent un nombre minuscule, comme
$10^{-17}$. On considère donc qu'un pivot est nul
quand sa valeur absolue est inférieure à un seuil
$\varepsilon$. Si le plus grand pivot possible de la
colonne $k$ est nul, la matrice est singulière.
\end{definition}

\subsection{Coût}

\begin{proposition}
Pour une matrice $n \times n$ :
\begin{itemize}
    \item l'élimination coûte environ
          $\frac{2}{3} n^3$ opérations ;
    \item la remontée coûte environ $n^2$ opérations.
\end{itemize}
\end{proposition}

\section{La décomposition LU}

L'élimination de Gauss « consomme » le second membre
$\vect{b}$. Pour résoudre le même système avec
beaucoup de seconds membres différents, on garde la
trace des opérations faites sur $A$.

\begin{theorem}[Décomposition $PA = LU$]
Pour toute matrice inversible $A \in \R^{n \times n}$,
il existe :
\begin{itemize}
    \item une \emph{matrice de permutation} $P$ : la
          matrice identité dont on a échangé des
          lignes, qui représente les échanges du
          pivot partiel ;
    \item une matrice $L$ triangulaire inférieure, avec
          des $1$ sur la diagonale, dont les
          coefficients sous la diagonale sont les
          multiplicateurs $\ell_{ik}$ ;
    \item une matrice $U$ triangulaire supérieure : le
          résultat de l'élimination de Gauss ;
\end{itemize}
telles que
\[
    PA = LU.
\]
\end{theorem}

\begin{methodbox}
Pour résoudre $A\vect{x} = \vect{b}$ une fois
$PA = LU$ connu :
\begin{enumerate}
    \item permuter le second membre : calculer
          $P\vect{b}$ ;
    \item \emph{descente} : résoudre
          $L\vect{y} = P\vect{b}$ de haut en bas
          ($L$ est triangulaire inférieure) ;
    \item \emph{remontée} : résoudre
          $U\vect{x} = \vect{y}$ de bas en haut.
\end{enumerate}
La décomposition coûte environ $\frac{2}{3} n^3$
opérations, mais elle est faite \textbf{une seule
fois}. Chaque résolution ne coûte ensuite qu'environ
$2n^2$ opérations.
\end{methodbox}

\begin{remark}
$L$ et $U$ tiennent dans \textbf{une seule} matrice
$n \times n$ : $U$ occupe la diagonale et le dessus,
et les multiplicateurs de $L$ occupent le dessous, à
la place des zéros créés par l'élimination. Les $1$
de la diagonale de $L$ ne sont pas stockés, puisqu'on
sait qu'ils valent $1$. De même, $P$ se représente par
un tableau de $n$ entiers : la ligne d'origine de
chaque ligne.
\end{remark}

\begin{importantbox}
Quand on échange deux lignes à la colonne $k$, les
multiplicateurs déjà calculés dans les colonnes
$1$ à $k - 1$ appartiennent à ces lignes : ils
doivent être échangés aussi. Sinon la solution peut
être juste, mais $PA \neq LU$.
\end{importantbox}

\section{Déterminant et inverse}

\subsection{Déterminant}

\begin{definition}[Déterminant d'une matrice $2 \times 2$]
\[
    \det \begin{pmatrix} a & b \\ c & d \end{pmatrix}
    = ad - bc.
\]
\end{definition}

\begin{proposition}
Pour $A, B \in \R^{n \times n}$ :
\begin{enumerate}
    \item $\det(AB) = \det(A) \det(B)$ ;
    \item le déterminant d'une matrice triangulaire est
          le produit de ses coefficients diagonaux ;
    \item échanger deux lignes multiplie le
          déterminant par $-1$ ;
    \item $A$ est inversible si et seulement si
          $\det(A) \neq 0$.
\end{enumerate}
\end{proposition}

\begin{proposition}[Déterminant par LU]
Si $PA = LU$ avec $s$ échanges de lignes effectifs,
alors
\[
    \det(A) = (-1)^s \prod_{i=1}^{n} u_{ii}.
\]
\end{proposition}

\begin{proof}
$\det(P)\det(A) = \det(L)\det(U)$. $L$ est
triangulaire avec des $1$ sur la diagonale, donc
$\det(L) = 1$. $U$ est triangulaire, donc
$\det(U) = \prod u_{ii}$. Enfin $P$ est l'identité
après $s$ échanges, donc $\det(P) = (-1)^s$, et
$\frac{1}{(-1)^s} = (-1)^s$.
\end{proof}

\subsection{Inverse}

\begin{methodbox}
La colonne $j$ de $A^{-1}$ est la solution de
$A\vect{x} = \vect{e}_j$, où $\vect{e}_j$ est la
colonne $j$ de $I_n$. Pour calculer $A^{-1}$ :
\begin{enumerate}
    \item décomposer $PA = LU$ \textbf{une fois} ;
    \item résoudre les $n$ systèmes
          $A\vect{x} = \vect{e}_j$ avec cette
          décomposition.
\end{enumerate}
Cela revient à résoudre $AX = I_n$, où le second
membre est une matrice : chaque colonne de $X$ est
un système à part.
\end{methodbox}

\section{Conditionnement}

\begin{definition}[Résidu]
Si $\tilde{\vect{x}}$ est la solution calculée, le
\emph{résidu} est $\vect{r} = \vect{b} - A\tilde{\vect{x}}$.
Un résidu de norme proche de $0$ signifie que
$\tilde{\vect{x}}$ vérifie presque le système.
\end{definition}

\begin{importantbox}
Un petit résidu ne garantit pas que
$\tilde{\vect{x}}$ est proche de la vraie solution.
Quand la matrice est \emph{mal conditionnée}, une
toute petite erreur sur $\vect{b}$ ou sur les arrondis
peut changer beaucoup la solution.
\end{importantbox}

\begin{definition}[Conditionnement]
Pour $A$ inversible, le \emph{conditionnement} est
\[
    \kappa(A) = \norm{A}\,\norm{A^{-1}} \ge 1.
\]
Plus $\kappa(A)$ est grand, plus la solution est
sensible aux erreurs. En \cpp|double|, on perd
environ $\log_{10} \kappa(A)$ chiffres significatifs
sur les 16 disponibles.
\end{definition}

\begin{example}[Matrice de Hilbert]
La matrice de Hilbert $H_n \in \R^{n \times n}$ est
définie par
\[
    h_{ij} = \frac{1}{i + j - 1}.
\]
Elle est inversible, mais son conditionnement grandit
très vite avec $n$. On choisit $\vect{x}$, on calcule
$\vect{b} = H_n \vect{x}$, on résout, et on compare la
solution obtenue au $\vect{x}$ choisi, pour
$n = 5$, $10$ et $15$.
\end{example}

\section{Exemples calculés à la main}

Les exercices de la note de l'étape 3 (le système
$3 \times 3$, le pivot nul et le pivot minuscule)
sont à faire à part, sur papier. Les exemples
ci-dessous utilisent d'autres matrices.

\begin{example}[Élimination et remontée]
\label[example]{ex-gauss}
On résout $A\vect{x} = \vect{b}$ avec
\[
    A = \begin{pmatrix} 2 & 1 & 1 \\ 4 & 3 & 3 \\ 8 & 7 & 9 \end{pmatrix},
    \qquad
    \vect{b} = \begin{pmatrix} 4 \\ 10 \\ 24 \end{pmatrix}.
\]
Colonne $1$, pivot $2$ : $\ell_{21} = 2$ et
$\ell_{31} = 4$.
\[
    \left(
    \begin{array}{ccc|c}
        2 & 1 & 1 & 4 \\
        4 & 3 & 3 & 10 \\
        8 & 7 & 9 & 24
    \end{array}
    \right)
    \xrightarrow[L_3 \leftarrow L_3 - 4L_1]{L_2 \leftarrow L_2 - 2L_1}
    \left(
    \begin{array}{ccc|c}
        2 & 1 & 1 & 4 \\
        0 & 1 & 1 & 2 \\
        0 & 3 & 5 & 8
    \end{array}
    \right).
\]
Colonne $2$, pivot $1$ : $\ell_{32} = 3$.
\[
    \xrightarrow{L_3 \leftarrow L_3 - 3L_2}
    \left(
    \begin{array}{ccc|c}
        2 & 1 & 1 & 4 \\
        0 & 1 & 1 & 2 \\
        0 & 0 & 2 & 2
    \end{array}
    \right).
\]
Remontée :
$x_3 = \frac{2}{2} = 1$, puis
$x_2 = 2 - x_3 = 1$, puis
$x_1 = \frac{4 - x_2 - x_3}{2} = 1$.
Donc $\vect{x} = (1, 1, 1)^\top$, et on vérifie que
$A\vect{x}$ redonne bien $(4, 10, 24)^\top$.

Ici, on a gardé le pivot diagonal pour simplifier.
Avec le pivot partiel, on aurait d'abord échangé
$L_1$ et $L_3$, car $\abs{8} > \abs{4} > \abs{2}$.
\end{example}

\begin{example}[Décomposition LU et déterminant]
Les multiplicateurs de l'\cref{ex-gauss} donnent,
sans échange de lignes ($P = I_3$) :
\[
    L = \begin{pmatrix} 1 & 0 & 0 \\ 2 & 1 & 0 \\ 4 & 3 & 1 \end{pmatrix},
    \qquad
    U = \begin{pmatrix} 2 & 1 & 1 \\ 0 & 1 & 1 \\ 0 & 0 & 2 \end{pmatrix}.
\]
On vérifie par exemple la ligne $3$ de $LU$ :
$4 \times (2, 1, 1) + 3 \times (0, 1, 1) + 1 \times (0, 0, 2)
= (8, 7, 9)$, c'est bien la ligne $3$ de $A$. Rangées
dans une seule matrice, $L$ et $U$ donnent
\[
    \begin{pmatrix} 2 & 1 & 1 \\ 2 & 1 & 1 \\ 4 & 3 & 2 \end{pmatrix}.
\]
Le déterminant vaut $\det(A) = 2 \times 1 \times 2 = 4$.
\end{example}

\begin{example}[Échange de lignes et signe du déterminant]
Pour
$A = \begin{pmatrix} 0 & 2 \\ 3 & 1 \end{pmatrix}$,
le pivot de la colonne $1$ est nul. Le pivot partiel
échange les deux lignes ($s = 1$) :
\[
    P = \begin{pmatrix} 0 & 1 \\ 1 & 0 \end{pmatrix},
    \qquad
    PA = \begin{pmatrix} 3 & 1 \\ 0 & 2 \end{pmatrix}
    = \underbrace{I_2}_{L}\ \underbrace{\begin{pmatrix} 3 & 1 \\ 0 & 2 \end{pmatrix}}_{U}.
\]
Donc $\det(A) = (-1)^1 \times 3 \times 2 = -6$, ce que
confirme la formule $2 \times 2$ :
$0 \times 1 - 2 \times 3 = -6$.
\end{example}

\begin{example}[Inverse d'une matrice $2 \times 2$]
Pour $A = \begin{pmatrix} 2 & 1 \\ 1 & 1 \end{pmatrix}$,
on résout $A\vect{x} = \vect{e}_1$ et
$A\vect{x} = \vect{e}_2$ :
\[
    \begin{cases} 2x + y = 1 \\ x + y = 0 \end{cases}
    \Rightarrow \begin{pmatrix} 1 \\ -1 \end{pmatrix},
    \qquad
    \begin{cases} 2x + y = 0 \\ x + y = 1 \end{cases}
    \Rightarrow \begin{pmatrix} -1 \\ 2 \end{pmatrix}.
\]
Ces deux solutions sont les colonnes de
$A^{-1} = \begin{pmatrix} 1 & -1 \\ -1 & 2 \end{pmatrix}$,
et on vérifie que $A A^{-1} = I_2$.
\end{example}

\begin{example}[Matrice singulière]
Pour $A = \begin{pmatrix} 1 & 2 \\ 2 & 4 \end{pmatrix}$,
la ligne $2$ vaut deux fois la ligne $1$. Après
$L_2 \leftarrow L_2 - 2L_1$, la ligne $2$ de $U$ est
nulle : le pivot de la colonne $2$ vaut $0$ et
$\det(A) = 0$.
\begin{itemize}
    \item Avec $\vect{b} = (3, 6)^\top$, la seconde
          équation devient $0 = 0$ : il y a une
          infinité de solutions.
    \item Avec $\vect{b} = (3, 7)^\top$, elle devient
          $0 = 1$ : il n'y a aucune solution.
\end{itemize}
Dans les deux cas, il n'y a pas de solution
\emph{unique}.
\end{example}

\section{Utilité dans le projet}

Un réseau de neurones n'inverse jamais de matrice, et
NNFS n'utilise pas ce chapitre directement. Il sert à
deux étapes de MiniNumpy :
\begin{itemize}
    \item les \textbf{moindres carrés} (étape 4) : la
          régression linéaire revient à résoudre les
          équations normales $A^\top A \vect{x} = A^\top \vect{b}$,
          un système linéaire ;
    \item la \textbf{méthode de Newton} (étape 6), qui
          résout un système linéaire à chaque
          itération.
\end{itemize}
Le conditionnement explique aussi un phénomène qu'on
reverra à l'entraînement : un calcul peut être
mathématiquement juste et numériquement faux.

\section{Implémentation}

\begin{codebox}
L'implémentation est l'objet de l'étape 3. Comme
pour le produit matriciel
(\cref{sec-matrices-implementation}), les
algorithmes de ce chapitre utiliseront la classe
\cpp|Matrix| sans en faire partie : ils seront
rangés dans \texttt{mnp/linalg.hpp} et
\texttt{src/linalg.cpp}, découpés en petites fonctions
(recherche du pivot, échange de lignes, élimination,
descente, remontée). Le second membre sera une
\cpp|Matrix| $n \times m$, ce qui permet de résoudre
plusieurs systèmes à la fois, et donc de calculer
l'inverse en résolvant $AX = I_n$. Le traitement d'une
matrice singulière (\cpp|std::optional| ou
exception) reste à décider.
\end{codebox}

\newpage
